% Generated draft: requires mathematical and source review.
\documentclass[11pt]{article}
\usepackage[margin=25mm]{geometry}
\usepackage{fontspec}
\setmainfont{Latin Modern Roman}
\usepackage{amsmath,amssymb,mathtools,unicode-math}
\setmathfont{Latin Modern Math}
\usepackage{enumitem,needspace,xurl,hyperref}
\setlength{\parindent}{0pt}
\setlength{\parskip}{5pt}
\setlength{\emergencystretch}{3em}
\newcommand{\N}{\mathbb{N}}
\newcommand{\Z}{\mathbb{Z}}
\newcommand{\Q}{\mathbb{Q}}
\newcommand{\R}{\mathbb{R}}
\newcommand{\C}{\mathbb{C}}
\newcommand{\F}{\mathbb{F}}
\newcommand{\A}{\mathbb{A}}
\newcommand{\PP}{\mathbb{P}}
\newcommand{\E}{\mathbb{E}}
\newcommand{\eps}{\varepsilon}
\newcommand{\dd}{\,\mathrm{d}}
\begin{document}
\section*{UnsolvedMath ID: 20000424}
\textbf{Source problem code:} catalogue entry
\subsection{Definitions and mathematical statement}
The question is attributed to K. Murty. Its source asks whether, for a fixed variety \(X/K\), the order of the pole of its \(L\)-function at the edge of the critical strip remains bounded when the ground field is enlarged, without assuming the Tate conjecture. The record does not specify the class of fields, the cohomological factor, or the permitted extensions. The following makes explicit the usual global-field interpretation; those choices require confirmation against the complete source.

Let \(K\) be a global field, let \(X/K\) be smooth and projective, fix \(0\leq m\leq\dim X\), and choose a prime \(\ell\ne\operatorname{char}K\). Put
\[
W=H^{2m}_{\mathrm{\acute et}}(X_{\overline K},\Q_\ell)(m),\qquad n=\dim_{\Q_\ell}W.
\]
For every finite separable extension \(L/K\) inside \(\overline K\), write \(\Gamma_L=\operatorname{Gal}(\overline K/L)\). With geometric Frobenius and a finite set \(S_L\) of excluded places containing the places needed to define the incomplete cohomological Euler product, write
\[
L^{S_L}_L(W,s)=\prod_{v\notin S_L}
\det\bigl(1-\operatorname{Frob}_v(Nv)^{-s}\mid W^{I_v}\bigr)^{-1}.
\]
The twist places the relevant edge at \(s=1\), corresponding to \(s=m+1\) for the untwisted \(H^{2m}\)-factor. Whenever this function has a meromorphic continuation to a neighborhood of \(1\), define its pole order there by
\[
p_L(W)=\max\{0,-\operatorname{ord}_{s=1}L^{S_L}_L(W,s)\}.
\]
The boundedness question in this interpretation is whether there is a finite constant \(C=C(X,K,m)\), independent of \(L\), such that
\[
p_L(W)\leq C\quad\text{for every finite separable }L/K,
\]
without assuming the Tate conjecture. The analytic continuation needed to speak of these orders must be established or explicitly assumed; it is not supplied by the wording of the question. Neither a uniform bound for all varieties nor uniformity in \(m\) is asserted.

The supplied research aid proposes the sharper bound \(C=n\) over global function fields, through the formula
\[
\dim W^{\Gamma_L}\ \leq\ -\operatorname{ord}_{s=1}L^{S_L}_L(W,s)
=[\mathbf 1:(W|_{\Gamma_L})^{\mathrm{ss}}]\ \leq\ n.
\]
Here \((\cdot)^{\mathrm{ss}}\) denotes semisimplification and \([\mathbf 1:V]\) the multiplicity of the trivial representation as a composition factor of \(V\). The aid attributes this to a corrected argument of Lyons using Lafforgue's work. This is a supplied, unreviewed partial-result claim, not an additional hypothesis of the original question or an independently verified theorem in this draft. In particular, equality of the middle quantity with \(\dim W^{\Gamma_L}\) cannot be inferred from this formula without an appropriate semisimplicity statement. The aid also reports a conditional automorphic route and a finite-monodromy case via Brauer induction and Hecke \(L\)-function nonvanishing, without giving their full hypotheses here.

The source's comment says that the boundedness is provable over function fields of finite fields and invokes a statement labeled \((\mathrm{T3})\) identifying pole order and Tate-class rank. It then starts a claim for abelian varieties ``in characteristic'' and stops. The missing characteristic and qualifications cannot be recovered from this record.
\subsection{Short English statement}
For a fixed cohomological part of a variety, ask for a bound on the edge-pole order that is independent of finite ground-field extension and does not assume the Tate conjecture. The usual global-field formulation is made explicit, while the truncated abelian-variety comment and the distinction between invariants and semisimplified constituents remain flagged.
\subsection{Sources}
\begin{itemize}
\item K. Murty, question in the AIM workshop problem list \emph{The Tate conjecture}, \url{https://aimath.org/WWN/tateconjecture/tateconjecture.pdf}. The extracted item is numbered 3; the supplied verification note says the official list numbers it 13, but does not supply the promised restored text.
\item Local registry \texttt{data/unsolvedmath/problems.json}, permanent UnsolvedMath ID \(20000424\), exact source code \texttt{catalogue entry}; canonical location \texttt{aim-algebraic-number-theory-notes.json}, notes[75].
\item The original comment names Lafforgue and \((\mathrm{T3})\). The attached research summary additionally names Lyons, Brauer induction, and Hecke \(L\)-function nonvanishing, but supplies no complete bibliographic citations or proof in this record. Research attempt 1 is explicitly machine-generated and not peer reviewed; it requires independent verification before citation.
\end{itemize}
\end{document}
